\documentclass[11pt,a4paper]{article}
\usepackage[T1]{fontenc}
\usepackage{lmodern,amsmath,amssymb,amsthm,mathtools,microtype}
\usepackage[margin=26mm]{geometry}
\usepackage[colorlinks=true,linkcolor=blue,citecolor=blue,urlcolor=blue]{hyperref}
\newtheorem{theorem}{Theorem}[section]
\newtheorem{proposition}[theorem]{Proposition}
\newtheorem{lemma}[theorem]{Lemma}
\newtheorem{corollary}[theorem]{Corollary}
\theoremstyle{definition}\newtheorem{remark}[theorem]{Remark}
\newcommand{\Q}{\mathbb Q}\newcommand{\Z}{\mathbb Z}
\newcommand{\hF}{h_{\mathrm F}}
\newcommand{\ord}{\operatorname{ord}}
\newcommand{\rad}{\operatorname{rad}}
\setlength{\emergencystretch}{2em}
\title{Scalar Characters and Contact-Height Estimates}
\author{}\date{10 November 2026}
\begin{document}\maketitle
\begin{abstract}
For the normalized quaternionic representation used in the five-contact
construction, the square of the nebentypus is the inverse square of the
scalar Dirichlet character. Choosing the local scalar characters of
minimal two-power order bounds the nebentypus by the largest two-part of
$p-1$ among the relevant primes. This gives a sharper modular degree
bound, an explicit cubic logarithmic contact estimate with constant less
than $2^{330670}$, and a finer estimate sensitive to congruence restrictions
on the contact primes. Separately, evaluation of differentials at a
primitive sixth root of unity removes the horizontal contribution to
the discriminant estimate for modular Jacobians. The arithmetic
realization remains an explicitly stated input.
\end{abstract}

\section{Scalar characters of minimal order}

For a positive squarefree integer $R$ divisible by $6$, define
\[
 U(R)=\max\bigl(\{2\}\cup
       \{2^{v_2(p-1)}:p\mid R,\ p\text{ odd prime}\}\bigr),
 \qquad E(R)=\prod_{p\mid R}(1+p^{-1}).
\]
Here $v_2(n)$ is the exponent of $2$ in the nonzero integer $n$.
Both $U(R)$ and $E(R)$ are effectively computable.

\begin{lemma}\label{lem:characters}
Prescribe a subset of the finite places dividing $R$ at which the
local boundary of a Dirichlet character under the squaring sequence is
to be nontrivial. The character used to realize these prescribed
boundaries can be chosen tame at every odd prime, of conductor at most
$4$ at $2$, and of order dividing $U(R)$. Moreover $U(R)\mid\varphi(R)$.
\end{lemma}
\begin{proof}
At an odd prime $p$, put $p-1=2^a u$ with $u$ odd. If a nontrivial
local boundary is prescribed, choose a generator $g$ of
$(\Z/p\Z)^\times$ and set $\chi_p(g)=\zeta_{2^a}$, a primitive
$2^a$-th root of unity. Then
\[
 \chi_p(-1)=\zeta_{2^a}^{(p-1)/2}=-1.
\]
Thus $\chi_p$ is odd and has the required boundary. The boundary
criterion is the local reciprocity calculation used in the scalar
lifting construction: the obstruction to taking a square root of a
finite local character is its value at $-1$. If the boundary is
trivial, use the trivial character. At $2$, use the character modulo
$4$ when needed. The product of these characters has order dividing
the least common multiple of their orders. All orders are powers of
$2$, so that least common multiple divides $U(R)$. The global boundary
is the prescribed sign class whenever the local prescription comes
from that class, by the same reciprocity and global-to-local argument
as in the scalar lifting construction. Finally each $2^{v_2(p-1)}$
divides $\varphi(R)=\prod_{p\mid R,\ p>2}(p-1)$, and this product is
even. The largest of these powers, together with $2$, therefore divides
$\varphi(R)$.
\end{proof}

\section{The determinant identity}

Let $A$ be an abelian surface over a finite normal extension $L/\Q$.
Suppose that compatible isogenies
$\mu_\sigma:{}^\sigma A\to A$ have been chosen, and that the normalized
semilinear operators
\begin{equation}\label{eq:normalized}
 \rho_\ell(\sigma)=b(\sigma)^{-1}a_\sigma\mu_\sigma\circ\sigma,
 \qquad a_\sigma=(\deg\mu_\sigma)^{-1/4},\qquad b^2=\chi,
\end{equation}
form a four-dimensional $\ell$-adic representation. Coefficient fields
are enlarged to contain these scalars, with the trivial Galois action
on the added coefficients. Assume that the quaternionic action gives
$\rho_\ell=r_\ell\oplus r_\ell$ and that modularity identifies
$r_\ell$ with a weight-two newform representation, in the Tate-module
convention
\[
                         \det r_\ell=\varepsilon\chi_{\rm cyc}.
\]
Here $\chi_{\rm cyc}$ denotes the $\ell$-adic cyclotomic character and
$\varepsilon$ is the nebentypus. These are precisely the normalized
representations and conventions in the modular construction of
\cite[proof of \texttt{o8-height}, (LR2)--(LR5)]{Source}.

\begin{proposition}\label{prop:determinant}
The character $\varepsilon\chi$ has order at most two. Consequently,
if $\chi$ is chosen as in Lemma~\ref{lem:characters}, then
\[
                         \ord(\varepsilon)\mid U(R).
\]
\end{proposition}
\begin{proof}
For an abelian surface, cup product and the trace isomorphism give
the natural identification
\[
                   \bigwedge^4 V_\ell(A)\simeq\Q_\ell(2).
\]
An isogeny acts on this determinant line by its degree. This follows
equally from its action on the top cohomology of a complex abelian
surface and comparison; the exterior-algebra description of abelian
cohomology is recalled in \cite[Chapter~I, Section~12]{Milne}.
Galois acts on the determinant line by $\chi_{\rm cyc}^2$.
Consequently, even though the isogeny has a conjugate source,
\[
 \det(\mu_\sigma\circ\sigma)
      =\deg(\mu_\sigma)\chi_{\rm cyc}(\sigma)^2.
\]
Taking determinants in~\eqref{eq:normalized}, and using
$a_\sigma^4\deg(\mu_\sigma)=1$, gives
\[
 \det\rho_\ell=b^{-4}\chi_{\rm cyc}^2
                 =\chi^{-2}\chi_{\rm cyc}^2.
\]
On the other hand,
$\det\rho_\ell=(\det r_\ell)^2
 =\varepsilon^2\chi_{\rm cyc}^2$.
Cancelling the cyclotomic factor proves
$\varepsilon^2=\chi^{-2}$. Thus
$\kappa=\varepsilon\chi$ is quadratic or trivial.
The orders of both $\chi$ and $\kappa$ divide $U(R)$, since $U(R)$
is even. Therefore the order of $\varepsilon=\kappa\chi^{-1}$
also divides $U(R)$.
\end{proof}

The quadratic twist of $b$ used to remove unwanted ramification
does not change this calculation: its fourth power is one, and its
square leaves $\chi$ unchanged. The proposition bounds the
nebentypus directly; it does not bound the newform coefficient-field
degree.

\section{A unit fibre for differential evaluation}

We use the original stable Faltings height. Let $X_H=X_1(N)/H$, where
$H\subset(\Z/N\Z)^\times$ contains $-1$, and let $J_H$ be its
Jacobian of positive dimension $g_H$. Write
\[
 d_H=[X_H:\mathbb P^1_j]
       =N[(\Z/N\Z)^\times:H]\prod_{p\mid N}(1+p^{-1}).
\]

\begin{proposition}\label{prop:unitfibre}
One has
\[
 \hF(J_H)\le\frac{d_H}{2}
           (\log6+\log N+2\log d_H)-\frac{g_H}{2}\log\pi.
\]
\end{proposition}
\begin{proof}
Use $t=j/1728$ and let $\zeta$ be a primitive sixth root of unity.
Choose a number field $K$ over which the curve is semistable, the
fibre $t=\zeta$ splits, and $\zeta$ is rational. Normalize
$\mathbb P^1_{\mathcal O_K}$ in $K(X_H)$ and specialize its affine
algebra at $t=\zeta$ to obtain an order $C\subset\mathcal O_K^{d_H}$.
The integral evaluation argument of \cite[Sections~1--3]{Evaluation}
maps the Hodge lattice into the trace-dual order. After projection to
$g_H$ coordinates on which evaluation is injective, the target's
arithmetic degree is at most half the logarithmic discriminant of $C$.
Injectivity holds because $d_H>2g_H-2$ for this Belyi map.

Both $\zeta$ and $1-\zeta$ are algebraic units. The fibre therefore
meets none of the horizontal branch divisors $0,1,\infty$ at finite
places. At primes greater than three, the torsion-different argument
of \cite[Lemma~4.2]{Evaluation} bounds the normalized vertical
contribution by $d_H\log N$. At $2$ and $3$, Lenstra's different
estimate gives at most $d_H(\log6+2\log d_H)$, exactly as in
\cite[Section~4]{Evaluation}. These arguments are unchanged after
the extension to $K$; its degree cancels after normalization.
Thus
\[
 \frac{\log\mathrm N\mathfrak D(C)}{[K:\Q]}
       \le d_H(\log6+\log N+2\log d_H).
\]

At every complex embedding, the distance from the image of $\zeta$
to each of $0$ and $1$ is one. Every disk of radius $s<1$ about
that image is unramified, and its inverse image consists of disjoint
analytic sheets. The mean-value inequality bounds evaluation of
$\omega/dt$ by $1/(\sqrt\pi s)$ times the $L^2$ norm of $\omega$.
Letting $s$ tend to one gives operator norm at most $1/\sqrt\pi$.
Hadamard's inequality bounds the determinant norm by $\pi^{-g_H/2}$.
The arithmetic degree inequality now gives the displayed result.
\end{proof}

\section{The refined contact estimate}

Put
\[
 B_* =\prod_{p\le76176}p,\qquad C_0=210^{300},\qquad
 \mathcal D=\{0,1,2401,-343/32\}.
\]
For $a=u/v\in\Q\setminus\mathcal D$ in lowest terms with $v>0$,
let $M=M_5(a)$ be the product of primes greater than $76176$ at
which at least one of
\[
 v,\quad u,\quad u-v,\quad u-2401v,\quad32u+343v
\]
has positive odd valuation. Write $h(a)=\log\max(|u|,v)$.

The arithmetic premise in this section is the normalized modular
construction of \cite[proof of \texttt{o8-height}]{Source}, including
its geometric inequality $h(a)\le2\max(0,\hF(A))+128$, its CM bound
$\hF(A)<854R$, and in the non-CM case its surface factor on a modular
Jacobian over $L'$, its level bound $N\le C_0R^2$, and
$\rad(N)\mid R$, where $R=B_*M$. The scalar characters in that
construction may be chosen as in Lemma~\ref{lem:characters}; its
conductor estimates use only tameness at odd primes and conductor at
most $4$ at $2$. Its field construction gives
$[L':\Q]\le1152\ord(\chi)$. This retains the source's arithmetic
premises and uses more information than a bare height inequality.

\begin{theorem}\label{thm:refined}
Under this premise, define
\[
 C_3'=19C_0B_*^3\log B_* ,\qquad
 C_2'=19C_0B_*^2E(B_*)\log B_* .
\]
Then
\begin{align}
 h(a)&\le C_3'M^3(1+\log M),\label{eq:uniform}\\
 h(a)&\le C_2'M^2 U(B_*M)E(M)(1+\log M).\label{eq:finer}
\end{align}
Moreover $C_3'<2^{330670}$, and hence
\[
                     h(a)<2^{330671}M^{7/2}.
\]
\end{theorem}
\begin{proof}
Let $H=\ker\varepsilon$. Proposition~\ref{prop:determinant} gives
\begin{equation}\label{eq:degree}
 d_H\le C_0R^2U(R)E(R)=:\mathcal T.
\end{equation}
Since $U(R)\mid\varphi(R)$,
\[
 U(R)E(R)\le\varphi(R)E(R)
        =R\prod_{p\mid R}(1-p^{-2})<R,
 \qquad \mathcal T<C_0R^3.
\]
Also $[L':\Q]\le1152U(R)\le1152R$.

For completeness, the numerical height transfer can be made with
the same estimates for both the refined and uniform bounds.
Write $\ell=\log R$. The exact inequalities $C_0^2<B_*$,
$144<B_*$ and $\ell\ge1$ give
\[
 \log N\le\tfrac52\ell,\qquad
 \log d_H\le\tfrac72\ell,\qquad
 \log12\le\tfrac12\ell.
\]
The weaker version of Proposition~\ref{prop:unitfibre} with
$\log12$ in place of $\log6$ and with additive term
$g_H\log(2/\sqrt\pi)$ yields
$\hF(J_H)\le6d_H\ell$.
Thus $H_J=\max(1,\hF(J_H)+3g_H/2)\le8\mathcal T\ell$.

R\'emond's subvariety inequality and the dimension-two isogeny
bound, in the form proved in \cite[Lemma~5.1]{Evaluation}, imply
\[
 \hF(A)\le H_J+8\left(32\log14+\log d'
           +\log\max(1,\log d',H_J)\right),\qquad d'=[L':\Q].
\]
Here $\log d'<9\ell$ and
$\log(8\mathcal T\ell)\le9\ell$, whereas
$8\mathcal T\ell>\max(1,\log d')$.
Using $\log14<3$ gives
\[
 \hF(A)\le8\mathcal T\ell+912\ell\le9\mathcal T\ell.
\]
The CM alternative is smaller than this bound as well.
Since $\mathcal T\ell>128$, the geometric comparison gives
$h(a)\le19\mathcal T\ell$. Substituting~\eqref{eq:degree}, using
$E(R)=E(B_*)E(M)$ and
$\log R\le(\log B_*)(1+\log M)$, proves~\eqref{eq:finer}.
The inequality $\mathcal T<C_0R^3$ proves~\eqref{eq:uniform}.

The exact checker described below verifies
\[
 19C_0B_*^3\cdot109445<2^{330670},\qquad B_*<2^{109445}.
\]
Since $\log2<1$, these inequalities imply $C_3'<2^{330670}$.
Finally $(1+u)e^{-u/2}<2$ for $u\ge0$, so the stated pure-power
bound follows on taking $u=\log M$.
\end{proof}

\begin{corollary}\label{cor:restricted}
Fix an integer $k\ge1$. Suppose that every prime dividing $M_5(a)$
satisfies $v_2(p-1)\le k$. Then
\[
 h(a)\le C_2'\max(65536,2^k)
                  M_5(a)^2(1+\log M_5(a))^2.
\]
In particular, for every fixed $\eta>0$ there is an effective constant
$C_{k,\eta}$ such that $h(a)\le C_{k,\eta}M_5(a)^{2+\eta}$
on this restricted set.
\end{corollary}
\begin{proof}
The primes in $B_*$ have largest two-part $65536$, attained at $65537$;
this is certified by exact trial division. Thus
$U(B_*M)\le\max(65536,2^k)$.
For squarefree $M$,
\[
 E(M)=\sum_{d\mid M}\frac1d\le\sum_{d=1}^{M}\frac1d
                                      \le1+\log M.
\]
Apply~\eqref{eq:finer}. The remaining logarithmic factor is absorbed
effectively by $M^\eta$, since
$\sup_{u\ge0}(1+u)^2e^{-\eta u}<\infty$.
\end{proof}

For example, primes $p\equiv17\pmod{32}$ satisfy $v_2(p-1)=4$.
The corollary applies when \emph{all} prime factors in the relevant
contact radical obey that restriction. The freedom to choose such
primes in a completeness argument does not impose the restriction on
arbitrary rational solutions in a soundness argument. No such
restriction on all contact primes is inferred here.

\begin{proposition}\label{prop:quadratic-obstruction}
For $\beta(X)=51X^2+2504$, the numbers
$U(B_*M_5(\beta(X)))$, with $X\in\Z$, are unbounded.
More precisely, for every $k\ge1$ there are infinitely many primes
$p>76176$ with $v_2(p-1)\ge k$ for which infinitely many integers
$X$ satisfy $v_p(\beta(X))=1$.
\end{proposition}
\begin{proof}
Set $L=2^{k+3}\cdot51\cdot2504$. Dirichlet's theorem on primes
in arithmetic progressions \cite[Section~1]{Kedlaya} supplies infinitely many primes
$p\equiv1\pmod L$. Discard the finitely many primes at most $76176$.
For every odd prime $q$ dividing $51\cdot2504$, quadratic
reciprocity gives $(q/p)=(p/q)=1$, because $p\equiv1\pmod{4q}$.
Also $(-1/p)=(2/p)=1$. Hence $-2504/51$ is a nonzero square
modulo $p$. Choose $x_0$ with $51x_0^2+2504\equiv0\pmod p$.
Its derivative $102x_0$ is nonzero modulo $p$.
Among the $p$ residue classes $x_0+pt\pmod{p^2}$, exactly one
makes $\beta$ vanish modulo $p^2$; any of the others has
$v_p(\beta)=1$. Each such residue class contains infinitely many
integers. For these inputs the contact at zero contributes $p$ to
$M_5(\beta(X))$, so $U(B_*M_5(\beta(X)))\ge2^{v_2(p-1)}\ge2^k$.
\end{proof}

Thus the congruence condition in Corollary~\ref{cor:restricted} cannot
hold with one fixed $k$ throughout this scalar family, even if the
input is already restricted to the rational integers.

\section{Exact verification and arithmetic scope}

The program
\begin{center}\small
\nolinkurl{computation/contact_heights/check_nebentypus_height.py}
\end{center}
constructs $B_*$ by a sieve and independently by trial division,
checks $U(B_*)=65536$, verifies the integer inequalities in
Theorem~\ref{thm:refined}, and tests the character-order identities in
cyclic groups. It also checks the two equations for a primitive sixth
root, $\zeta^2-\zeta+1=0$ and $\zeta(1-\zeta)=1$.
These finite checks verify the constants and elementary algebra.
The determinant, duality, and arithmetic-realization assertions are
the mathematical proofs and premises stated above, not conclusions
of the computation.

\begin{thebibliography}{9}\small
\bibitem{Source}
E.~Glazer and contributors, \emph{Diophantine definition of the integers
in the rationals}, \texttt{manuscripts/quartic\_v9.tex}, commit
\texttt{14e8f601fe0d9173d4a5f084aedc3e07066a7688},
\url{https://github.com/elliotglazer/diophantine-z-q}.
\bibitem{Evaluation}
J.~Brox and contributors, \emph{Evaluation of Differentials and Heights
of Modular Jacobians},
\texttt{evaluation\_of\_differentials\_20261109\_v17.tex},
\url{https://github.com/JoseBrox/diophantine-z-q_improvements}.
\bibitem{Milne}
J.~S. Milne, \emph{Abelian Varieties}, course notes, version 2.00,
2008, Chapter~I, Section~12,
\url{https://www.jmilne.org/math/CourseNotes/AV.pdf}.
\bibitem{Kedlaya}
K.~S. Kedlaya, \emph{Primes in arithmetic progressions},
MIT 18.785 lecture notes, Spring 2007, Section~1,
\url{https://kskedlaya.org/18.785/dirichlet.pdf}.
\bibitem{Javanpeykar}
A.~Javanpeykar, with an appendix by P.~Bruin,
\emph{Polynomial bounds for Arakelov invariants of Belyi curves},
Algebra \& Number Theory \textbf{8} (2014), 89--140,
\url{https://doi.org/10.2140/ant.2014.8.89}.
\bibitem{Remond}
G.~R\'emond, \emph{Propri\'et\'es de la hauteur de Faltings},
Annali della Scuola Normale Superiore di Pisa, Classe di Scienze
\textbf{23} (2022), 1589--1596,
\url{https://doi.org/10.2422/2036-2145.202010_062}.
\bibitem{GaudronRemond}
\'E.~Gaudron and G.~R\'emond,
\emph{Nouveaux th\'eor\`emes d'isog\'enie},
M\'emoires de la Soci\'et\'e Math\'ematique de France
\textbf{176} (2023), Theorem~1.9(1),
\url{https://doi.org/10.24033/msmf.484}.
\end{thebibliography}
\end{document}
